\chapter{Testes Diagnósticos}

\section{Após uma regressão}

`test(model, ...)` e outros comandos de diagnóstico verificam se os
pressupostos por trás de um estimador são plausíveis. Este capítulo
cobre os testes pós-estimação e autônomos mais comuns.

\section{Heterocedasticidade}

\begin{lstlisting}[caption={Testes de White e Breusch-Pagan}]
test(m, "white")     // Teste de White
test(m, "bp")        // Breusch-Pagan
\end{lstlisting}

White é um teste geral baseado na regressão dos resíduos quadrados
sobre produtos cruzados e quadrados de $X$. Breusch-Pagan é um teste
multiplicador de Lagrange sob normalidade.

\section{Autocorrelação}

\begin{lstlisting}[caption={Testes de autocorrelação}]
test(m, "dw")        // Durbin-Watson
test(m, "bg")        // Breusch-Godfrey LM
cusumtest(m)         // CUSUM de estabilidade
\end{lstlisting}

Durbin-Watson testa autocorrelação de primeira ordem nos resíduos de
OLS. Breusch-Godfrey testa até um defasamento especificado.
`cusumtest` verifica se resíduos recursivos permanecem dentro dos
limites de 5\%.

\section{Normalidade}

\begin{lstlisting}[caption={Testes de normalidade}]
swilk(df, y)         // Shapiro-Wilk
sfrancia(df, y)      // Shapiro-Francia
sktest(df, y)        // Assimetria/curtose (Jarque-Bera + D'Agostino)
jb(m)                // Jarque-Bera nos resíduos
\end{lstlisting}

\section{Especificação}

\begin{lstlisting}[caption={RESET e link test}]
reset(m)             // Ramsey RESET
linktest(m)          // Link test para logit/probit
\end{lstlisting}

RESET adiciona potências dos valores ajustados à regressão; um
coeficiente significativo indica não-linearidade omitida. `linktest` é
o análogo para modelos GLM.

\section{Outliers e influência}

\begin{lstlisting}[caption={Diagnósticos de influência}]
influence(m)         // DFFITS, alavancagem
leverage(m)          // alavancagem
cooks(m)             // distância de Cook
\end{lstlisting}

Alta alavancagem significa que uma observação tem regressores pouco
comuns. Alto DFFITS ou distância de Cook significa que uma observação
altera os coeficientes ajustados.

\section{Multicolinearidade}

\begin{lstlisting}[caption={VIF}]
vif(m)               // fatores de inflação de variância
\end{lstlisting}

Um VIF acima de 10 (ou às vezes 5) sugere multicolinearidade
problemática.

\section{Raiz unitária}

\begin{lstlisting}[caption={Testes de raiz unitária}]
adf(df, y)           // Dickey-Fuller Aumentado
kpss(df, y)          // KPSS (nula de estacionariedade)
pp(df, y)            // Phillips-Perron
zivot(df, y)         // Zivot-Andrews (quebra estrutural)
\end{lstlisting}

ADF, PP e KPSS são os três testes clássicos. KPSS tem nula de
estacionariedade, complementando ADF e PP, que têm nula de raiz
unitária.

\section{Cointegração}

\begin{lstlisting}[caption={Testes de cointegração}]
engle_granger(df, y, x, lags=1)   // equação única
johansen(df, y, x, lags=1)        // multivariado
\end{lstlisting}

Engle-Granger é simples mas assume um único vetor de cointegração.
Johansen testa o número de vetores de cointegração em um sistema
multivariado.

\section{Testes não paramétricos}

\begin{lstlisting}[caption={Testes não paramétricos}]
spearman(df, x, y)        // Correlação de Spearman
ranksum(df, x, by=group)  // Wilcoxon rank-sum
kruskal(df, y, by=group)  // Kruskal-Wallis
signrank(df, x, y)        // Wilcoxon signed-rank
\end{lstlisting}

São alternativas livres de distribuição para testes t, ANOVA e
correlação de Pearson.
